% GENERATED FILE. Edit canonical lessons, then run npm run generate:books.
% canonical-source-hash: fnv1a64:48317dbdf1e81363
% canonical-lessons: ML-C194-prasavam, ML-C194-opparesan, ML-C194-kelvi, ML-C194-arogyam, ML-C194-kanpola

\chapter{Childbirth, Surgical operation, Hearing, Health, Eyelid}
\label{ch:ml-childbirth-surgical-operation-hearing-health-eyelid}
\hlchaptermodality{194}

\begin{chapteropening}
\textbf{By the end of this chapter:} \emph{I can say childbirth, a surgical operation, hearing, health and an eyelid.}

Complete the last lesson of chapter 194: I can say childbirth, a surgical operation, hearing, health and an eyelid.
\end{chapteropening}

\section[\texorpdfstring{\ml{പ്രസവം} (prasavaṁ)}{പ്രസവം (prasavaṁ)}]{\ml{പ്രസവം} (prasavaṁ) — childbirth, delivery}

\label{lesson:ML-C194-prasavam}



\begin{quote}
Before the new one: say the Malayalam for pride, then the Malayalam for wonder.
\end{quote}

\subsection*{You'll want to know: \ml{പ്രസവം}}
\textbf{\ml{പ്രസവം}} — \emph{prasavaṁ} — \textquotedblleft{}childbirth, delivery\textquotedblright{}.

\begin{cousinweb}
\textbf{\ml{പ്രസവം}} (\emph{prasavaṁ}) $\leftarrow$ Sanskrit \emph{prasava}, \textquotedblleft{}procreation, bringing forth\textquotedblright{}, taken over whole as a tatsama, with Malayalam's noun ending \emph{-ṁ} (CDIAL 8827).
\end{cousinweb}

\begin{grammarlens}[title={What you've built}]
One more word for this chapter.
\end{grammarlens}

\subsection*{Your turn}
\begin{itemize}
\raggedright
  \item \emph{Say it:} \emph{prasavaṁ}
  \item \emph{Say it:} \emph{prasavaṁ}, once more
  \item \emph{Say it:} say \emph{prasavaṁ}
  \item \emph{Recall:} say the Malayalam for pride, then the Malayalam for wonder, then say \emph{prasavaṁ} again
\end{itemize}

\begin{tcolorbox}[breakable,colback=teal!4,colframe=teal!35!black,title={Before you move on}]
What does \ml{പ്രസവം} mean? (\textquotedblleft{}Childbirth, delivery\textquotedblright{}.) Say it once more.
\end{tcolorbox}
\section[\texorpdfstring{\ml{ഓപ്പറേഷൻ} (ōppaṟēṣan)}{ഓപ്പറേഷൻ (ōppaṟēṣan)}]{\ml{ഓപ്പറേഷൻ} (ōppaṟēṣan) — a surgical operation}

\label{lesson:ML-C194-opparesan}



\begin{quote}
Before the new one: say the Malayalam for wonder, then the Malayalam for childbirth.
\end{quote}

\subsection*{You'll want to know: \ml{ഓപ്പറേഷൻ}}
\textbf{\ml{ഓപ്പറേഷൻ}} — \emph{ōppaṟēṣan} — \textquotedblleft{}a surgical operation\textquotedblright{}.

\begin{cousinweb}
\textbf{\ml{ഓപ്പറേഷൻ}} (\emph{ōppaṟēṣan}) $\leftarrow$ English \emph{operation}, from Latin \emph{operātiō}, \textquotedblleft{}a working\textquotedblright{}.
\end{cousinweb}

\begin{grammarlens}[title={What you've built}]
One more word for this chapter.
\end{grammarlens}

\subsection*{Your turn}
\begin{itemize}
\raggedright
  \item \emph{Say it:} \emph{ōppaṟēṣan}
  \item \emph{Say it:} \emph{ōppaṟēṣan}, once more
  \item \emph{Say it:} say \emph{ōppaṟēṣan}
  \item \emph{Recall:} say the Malayalam for wonder, then the Malayalam for childbirth, then say \emph{ōppaṟēṣan} again
\end{itemize}

\begin{tcolorbox}[breakable,colback=teal!4,colframe=teal!35!black,title={Before you move on}]
What does \ml{ഓപ്പറേഷൻ} mean? (\textquotedblleft{}A surgical operation\textquotedblright{}.) Say it once more.
\end{tcolorbox}
\section[\texorpdfstring{\ml{കേൾവി} (kēḷvi)}{കേൾവി (kēḷvi)}]{\ml{കേൾവി} (kēḷvi) — hearing}

\label{lesson:ML-C194-kelvi}



\begin{quote}
Before the new one: say the Malayalam for childbirth, then the Malayalam for a surgical operation.
\end{quote}

\subsection*{You'll want to know: \ml{കേൾവി}}
\textbf{\ml{കേൾവി}} — \emph{kēḷvi} — \textquotedblleft{}hearing\textquotedblright{}.

\begin{cousinweb}
\textbf{\ml{കേൾവി}} (\emph{kēḷvi}) is native Dravidian, the noun of the verb \textquotedblleft{}to hear\textquotedblright{}; DEDR 2017 lists Malayalam \emph{kēḷvi}, \textquotedblleft{}hearing\textquotedblright{}, itself, beside Tamil \emph{kēḷ} and Kannada \emph{kēḷ}.
\end{cousinweb}

\begin{grammarlens}[title={What you've built}]
One more word for this chapter.
\end{grammarlens}

\subsection*{Your turn}
\begin{itemize}
\raggedright
  \item \emph{Say it:} \emph{kēḷvi}
  \item \emph{Say it:} \emph{kēḷvi}, once more
  \item \emph{Say it:} say \emph{kēḷvi}
  \item \emph{Recall:} say the Malayalam for childbirth, then the Malayalam for a surgical operation, then say \emph{kēḷvi} again
\end{itemize}

\begin{tcolorbox}[breakable,colback=teal!4,colframe=teal!35!black,title={Before you move on}]
What does \ml{കേൾവി} mean? (\textquotedblleft{}Hearing\textquotedblright{}.) Say it once more.
\end{tcolorbox}
\section[\texorpdfstring{\ml{ആരോഗ്യം} (ārōgyaṁ)}{ആരോഗ്യം (ārōgyaṁ)}]{\ml{ആരോഗ്യം} (ārōgyaṁ) — health}

\label{lesson:ML-C194-arogyam}



\begin{quote}
Before the new one: say the Malayalam for a surgical operation, then the Malayalam for hearing.
\end{quote}

\subsection*{You'll want to know: \ml{ആരോഗ്യം}}
\textbf{\ml{ആരോഗ്യം}} — \emph{ārōgyaṁ} — \textquotedblleft{}health\textquotedblright{}.

\begin{cousinweb}
\textbf{\ml{ആരോഗ്യം}} (\emph{ārōgyaṁ}) $\leftarrow$ Sanskrit \emph{ārōgya}, \textquotedblleft{}health\textquotedblright{}, taken over whole as a tatsama, with Malayalam's noun ending \emph{-ṁ} (CDIAL 1330). Its parts are \emph{a-}, \textquotedblleft{}without\textquotedblright{}, and \emph{rōga}, \textquotedblleft{}disease\textquotedblright{}.
\end{cousinweb}

\begin{grammarlens}[title={What you've built}]
One more word for this chapter.
\end{grammarlens}

\subsection*{Your turn}
\begin{itemize}
\raggedright
  \item \emph{Say it:} \emph{ārōgyaṁ}
  \item \emph{Say it:} \emph{ārōgyaṁ}, once more
  \item \emph{Say it:} say \emph{ārōgyaṁ}
  \item \emph{Recall:} say the Malayalam for a surgical operation, then the Malayalam for hearing, then say \emph{ārōgyaṁ} again
\end{itemize}

\begin{tcolorbox}[breakable,colback=teal!4,colframe=teal!35!black,title={Before you move on}]
What does \ml{ആരോഗ്യം} mean? (\textquotedblleft{}Health\textquotedblright{}.) Say it once more.
\end{tcolorbox}
\section[\texorpdfstring{\ml{കൺപോള} (kaṇpōḷa)}{കൺപോള (kaṇpōḷa)}]{\ml{കൺപോള} (kaṇpōḷa) — an eyelid}

\label{lesson:ML-C194-kanpola}



\begin{quote}
Before the new one: say the Malayalam for hearing, then the Malayalam for health.
\end{quote}

\subsection*{You'll want to know: \ml{കൺപോള}}
\textbf{\ml{കൺപോള}} — \emph{kaṇpōḷa} — \textquotedblleft{}an eyelid\textquotedblright{}.

\begin{cousinweb}
\textbf{\ml{കൺപോള}} (\emph{kaṇpōḷa}) is native Dravidian: \emph{kaṇ}, \textquotedblleft{}eye\textquotedblright{} (DEDR 1159, with Tamil \emph{kaṇ}), and \emph{pōḷa}, which may be the word DEDR 4455 lists as \textquotedblleft{}bubble, blister\textquotedblright{}; that second link is not certain.
\end{cousinweb}

\begin{grammarlens}[title={What you've built}]
One more word for this chapter.
\end{grammarlens}

\subsection*{Your turn}
\begin{itemize}
\raggedright
  \item \emph{Say it:} \emph{kaṇpōḷa}
  \item \emph{Say it:} \emph{kaṇpōḷa}, once more
  \item \emph{Say it:} say \emph{kaṇpōḷa}
  \item \emph{Recall:} say the Malayalam for hearing, then the Malayalam for health, then say \emph{kaṇpōḷa} again
\end{itemize}

\begin{tcolorbox}[breakable,colback=teal!4,colframe=teal!35!black,title={Before you move on}]
What does \ml{കൺപോള} mean? (\textquotedblleft{}An eyelid\textquotedblright{}.) Say it once more.
\end{tcolorbox}
